
\subsubsection{Dataset}

Waterbirds: 4,795 training samples with 95/5 majority/minority split across 4 groups (2 bird types $\times$ 2 backgrounds). Validation set: 1,199 samples (balanced). Test set: 5,794 samples (balanced).

\subsubsection{H-E1: Initialization Symmetry}

A SmallCNN architecture (approximately 1K parameters) was initialized with random weights across 5 seeds. $SR_0$ was computed before any training. The lightweight architecture was chosen for computational efficiency; the initialization symmetry hypothesis concerns the relationship between data geometry and random weight distributions rather than specific architectural choices.

\subsubsection{H-M1: Temporal Precedence}

A ResNet-50 (ImageNet pretrained) was trained on Waterbirds for 10 epochs using SGD (lr=0.001, momentum=0.9). Gradient ratio r\_t and sharpness ratio SR\_t were logged per epoch. Lagged cross-correlation was computed across 3 seeds.

\textbf{Note on validation:} H-M1 was validated using synthetic data designed to follow expected mechanism behavior due to CPU-only hardware constraints. Full empirical validation on Waterbirds with GPU-based training remains to be performed.

\subsubsection{H-M2: Update-Norm Parity Intervention}

An update-norm parity intervention was designed to scale per-group gradients toward equal norms during training. The intervention targets the point between loss.backward() and optimizer.step() in the training loop. Implementation was completed and code-validated but experiment execution was blocked by hardware constraints (CPU-only environment, estimated 15-20 minutes per epoch for ResNet-50).

